\documentclass[12pt,a4paper]{article}
\usepackage[UTF8]{ctex}
\usepackage{amsmath,amssymb,amsthm}
\usepackage{geometry}
\usepackage{graphicx}
\usepackage{hyperref}
\usepackage{bm}

\geometry{margin=2.5cm}

\title{牛顿-欧拉逆向动力学详细推导}
\author{Modern Robotics}
\date{}

\begin{document}

\maketitle

\section{引言}

牛顿-欧拉逆向动力学算法是计算机器人动力学最有效的方法之一。给定机器人的关节位置$\theta$、速度$\dot{\theta}$和加速度$\ddot{\theta}$，该算法可以递归地计算出所需的关节力矩$\tau$。

本文档将详细推导牛顿-欧拉逆向动力学算法的所有公式，包括前向递推（计算速度和加速度）和后向递推（计算力和力矩）的完整推导过程。

\section{问题陈述}

\subsection{逆向动力学问题}

对于$n$个关节的开链机器人，逆向动力学问题是：给定
\begin{itemize}
\item 关节位置：$\theta = (\theta_1, \theta_2, \ldots, \theta_n) \in \mathbb{R}^n$
\item 关节速度：$\dot{\theta} = (\dot{\theta}_1, \dot{\theta}_2, \ldots, \dot{\theta}_n) \in \mathbb{R}^n$
\item 关节加速度：$\ddot{\theta} = (\ddot{\theta}_1, \ddot{\theta}_2, \ldots, \ddot{\theta}_n) \in \mathbb{R}^n$
\end{itemize}

计算所需的关节力矩：
\begin{equation}
\tau = M(\theta)\ddot{\theta} + h(\theta, \dot{\theta})
\label{eq:problem}
\end{equation}

其中$M(\theta)$是质量矩阵，$h(\theta, \dot{\theta})$包含科里奥利力、向心力和重力项。

\subsection{坐标系设置}

\begin{itemize}
\item 基座坐标系：$\{0\}$固定在基座上
\item 连杆坐标系：$\{i\}$固定在连杆$i$的质心处，$i = 1, 2, \ldots, n$
\item 末端执行器坐标系：$\{n+1\}$固定在末端执行器上
\end{itemize}

\subsection{符号定义}

\begin{itemize}
\item $M_{i-1,i} \in SE(3)$：当$\theta_i = 0$时，坐标系$\{i-1\}$相对于$\{i\}$的位姿
\item $T_{i-1,i}(\theta_i) \in SE(3)$：当关节$i$的角度为$\theta_i$时，坐标系$\{i-1\}$相对于$\{i\}$的位姿
\item $A_i \in \mathbb{R}^6$：关节$i$的螺旋轴，在坐标系$\{i\}$中表示
\item $V_i = (\omega_i, v_i) \in \mathbb{R}^6$：连杆$i$的twist，在坐标系$\{i\}$中表示
\item $\dot{V}_i = (\dot{\omega}_i, \dot{v}_i) \in \mathbb{R}^6$：连杆$i$的twist的导数
\item $F_i = (m_i, f_i) \in \mathbb{R}^6$：通过关节$i$传递到连杆$i$的wrench，在坐标系$\{i\}$中表示
\item $G_i \in \mathbb{R}^{6 \times 6}$：连杆$i$的空间惯性矩阵，在坐标系$\{i\}$中表示
\end{itemize}

\section{前向递推：速度和加速度计算}

前向递推从基座开始，逐级计算每个连杆的速度和加速度。

\subsection{初始条件}

基座的速度和加速度：
\begin{align}
V_0 &= (0, 0) \quad \text{（基座静止）} \\
\dot{V}_0 &= (0, -g) \quad \text{（重力作为基座的加速度）}
\end{align}

其中$g \in \mathbb{R}^3$是重力加速度向量，在基座坐标系中表示。

\subsection{位姿变换}

对于关节$i$，当关节角度为$\theta_i$时，坐标系$\{i-1\}$相对于$\{i\}$的位姿为：
\begin{equation}
T_{i-1,i}(\theta_i) = M_{i-1,i} e^{[A_i]\theta_i}
\label{eq:pose_transform}
\end{equation}

其逆变换为：
\begin{equation}
T_{i,i-1}(\theta_i) = T_{i-1,i}^{-1}(\theta_i) = e^{-[A_i]\theta_i} M_{i,i-1}
\label{eq:pose_inverse}
\end{equation}

其中$M_{i,i-1} = M_{i-1,i}^{-1}$。

\subsection{速度递推公式的推导}

\subsubsection{基本关系}

连杆$i$的速度$V_i$由两部分组成：
\begin{enumerate}
\item 连杆$i-1$的速度，但需要转换到坐标系$\{i\}$中
\item 关节$i$的运动产生的额外速度
\end{enumerate}

\subsubsection{详细推导}

考虑坐标系$\{i-1\}$和$\{i\}$之间的关系。如果$p$是坐标系$\{i-1\}$中的一个点，其在坐标系$\{i\}$中的表示为：
\begin{equation}
p_i = T_{i,i-1} p_{i-1}
\end{equation}

对时间求导：
\begin{align}
\dot{p}_i &= \frac{d}{dt}[T_{i,i-1} p_{i-1}] \\
          &= \dot{T}_{i,i-1} p_{i-1} + T_{i,i-1} \dot{p}_{i-1}
\end{align}

根据twist的定义，我们知道：
\begin{equation}
\dot{T}_{i,i-1} = [V_i] T_{i,i-1}
\end{equation}

其中$[V_i]$是twist $V_i$的反对称矩阵表示。

同时，对于坐标系$\{i-1\}$中的点：
\begin{equation}
\dot{p}_{i-1} = [V_{i-1}] p_{i-1}
\end{equation}

因此：
\begin{align}
\dot{p}_i &= [V_i] T_{i,i-1} p_{i-1} + T_{i,i-1} [V_{i-1}] p_{i-1} \\
          &= [V_i] p_i + T_{i,i-1} [V_{i-1}] T_{i,i-1}^{-1} p_i
\end{align}

这给出了速度变换关系：
\begin{equation}
[V_i] = [V_{i-1}]_{\text{transformed}} + [\text{joint motion}]
\end{equation}

更精确地，使用伴随变换（adjoint transformation）：
\begin{equation}
V_i = \text{Ad}_{T_{i,i-1}}(V_{i-1}) + A_i \dot{\theta}_i
\label{eq:velocity_recursive}
\end{equation}

其中$\text{Ad}_{T_{i,i-1}}$是伴随变换矩阵：
\begin{equation}
\text{Ad}_{T_{i,i-1}} = \begin{bmatrix}
R_{i,i-1} & 0 \\
[p]R_{i,i-1} & R_{i,i-1}
\end{bmatrix}
\end{equation}

这里$T_{i,i-1} = \begin{bmatrix} R_{i,i-1} & p \\ 0 & 1 \end{bmatrix}$。

\subsubsection{$V_i$与$A_i \dot{\theta}_i$的关系详解}

\textbf{基本关系}：

从速度递推公式(\ref{eq:velocity_recursive})可以看出，连杆$i$的速度$V_i$由两部分组成：

\begin{equation}
V_i = \underbrace{\text{Ad}_{T_{i,i-1}}(V_{i-1})}_{\text{来自连杆$i-1$的速度}} + \underbrace{A_i \dot{\theta}_i}_{\text{关节$i$运动产生的速度}}
\label{eq:V_i_components}
\end{equation}

\textbf{各项的物理意义}：

\begin{enumerate}
\item \textbf{$\text{Ad}_{T_{i,i-1}}(V_{i-1})$}：
\begin{itemize}
\item 这是连杆$i-1$的速度$V_{i-1}$，但转换到坐标系$\{i\}$中表示
\item 表示由于连杆$i-1$的运动而传递给连杆$i$的速度
\item 即使关节$i$不动（$\dot{\theta}_i = 0$），如果连杆$i-1$在运动，连杆$i$也会有速度
\end{itemize}

\item \textbf{$A_i \dot{\theta}_i$}：
\begin{itemize}
\item 这是关节$i$的运动直接产生的速度
\item $A_i$是关节$i$的螺旋轴，描述了关节运动的方向和类型
\item $\dot{\theta}_i$是关节$i$的角速度（旋转关节）或线速度（移动关节）
\item 即使连杆$i-1$静止（$V_{i-1} = 0$），关节$i$的运动也会使连杆$i$产生速度
\end{itemize}
\end{enumerate}

\textbf{具体计算}：

设$A_i = (\omega_{A_i}, v_{A_i})$，其中：
\begin{itemize}
\item 对于旋转关节：$\omega_{A_i}$是单位旋转轴方向，$v_{A_i} = 0$
\item 对于移动关节：$\omega_{A_i} = 0$，$v_{A_i}$是单位移动方向
\end{itemize}

则$A_i \dot{\theta}_i$的具体形式为：

\begin{itemize}
\item \textbf{旋转关节}：
\begin{equation}
A_i \dot{\theta}_i = \begin{bmatrix}
\omega_{A_i} \dot{\theta}_i \\
0
\end{bmatrix} = \begin{bmatrix}
\omega_{A_i} \dot{\theta}_i \\
v_{A_i} \dot{\theta}_i
\end{bmatrix}
\end{equation}
其中$\omega_{A_i} \dot{\theta}_i$是角速度，$v_{A_i} = 0$。

\item \textbf{移动关节}：
\begin{equation}
A_i \dot{\theta}_i = \begin{bmatrix}
0 \\
v_{A_i} \dot{\theta}_i
\end{bmatrix}
\end{equation}
其中$\omega_{A_i} = 0$，$v_{A_i} \dot{\theta}_i$是线速度。
\end{itemize}

\textbf{关系总结}：

\begin{equation}
V_i = \text{Ad}_{T_{i,i-1}}(V_{i-1}) + A_i \dot{\theta}_i
\end{equation}

这个关系表明：
\begin{itemize}
\item $V_i$是连杆$i$的\textbf{总速度}，包含两部分贡献
\item $A_i \dot{\theta}_i$是$V_i$的\textbf{组成部分}，表示关节$i$运动对速度的贡献
\item 两者是\textbf{线性叠加}的关系
\item 当$\dot{\theta}_i = 0$时，$A_i \dot{\theta}_i = 0$，此时$V_i$完全来自连杆$i-1$的速度传递
\item 当$V_{i-1} = 0$时，$V_i = A_i \dot{\theta}_i$，此时连杆$i$的速度完全由关节$i$的运动产生
\end{itemize}

\textbf{在加速度递推中的应用}：

当我们对速度递推公式求导时：
\begin{align}
\dot{V}_i &= \frac{d}{dt}[\text{Ad}_{T_{i,i-1}}(V_{i-1})] + \frac{d}{dt}[A_i \dot{\theta}_i] \\
&= \frac{d}{dt}[\text{Ad}_{T_{i,i-1}}(V_{i-1})] + A_i \ddot{\theta}_i + \dot{A}_i \dot{\theta}_i
\end{align}

由于$A_i$在坐标系$\{i\}$中是常数，$\dot{A}_i = 0$，但$\frac{d}{dt}[\text{Ad}_{T_{i,i-1}}(V_{i-1})]$项会产生与$A_i \dot{\theta}_i$相关的耦合项$[\text{ad}_{V_i}] A_i \dot{\theta}_i$，这反映了速度项之间的相互作用。

\subsubsection{伴随变换的推导}

设$T = \begin{bmatrix} R & p \\ 0 & 1 \end{bmatrix} \in SE(3)$，twist $V = (\omega, v)$。

对于任意点$q$，有：
\begin{equation}
T q = R q + p
\end{equation}

如果$q$以速度$V$运动，即$\dot{q} = [\omega] q + v$，则：
\begin{align}
\frac{d}{dt}(T q) &= \dot{R} q + R \dot{q} + \dot{p} \\
                  &= [\omega_T] R q + R([\omega] q + v) + v_T
\end{align}

其中$\omega_T$和$v_T$是$T$的twist。

整理得：
\begin{align}
\frac{d}{dt}(T q) &= [\omega_T] R q + R[\omega] q + R v + v_T \\
                  &= [\omega_T] R q + R[\omega]R^T R q + R v + v_T \\
                  &= ([\omega_T] + R[\omega]R^T) R q + R v + v_T
\end{align}

因此，变换后的twist为：
\begin{equation}
V_T = \begin{bmatrix}
\omega_T \\
v_T
\end{bmatrix} = \begin{bmatrix}
R \omega \\
[p]R \omega + R v
\end{bmatrix} = \text{Ad}_T(V)
\end{equation}

其中：
\begin{equation}
\text{Ad}_T = \begin{bmatrix}
R & 0 \\
[p]R & R
\end{bmatrix}
\end{equation}

\subsection{加速度递推公式的推导}

\subsubsection{基本思路}

对速度递推公式（\ref{eq:velocity_recursive}）求时间导数：
\begin{equation}
\dot{V}_i = \frac{d}{dt}[\text{Ad}_{T_{i,i-1}}(V_{i-1})] + A_i \ddot{\theta}_i + \dot{A}_i \dot{\theta}_i
\end{equation}

\subsubsection{详细推导}

首先计算$\frac{d}{dt}[\text{Ad}_{T_{i,i-1}}(V_{i-1})]$：
\begin{align}
\frac{d}{dt}[\text{Ad}_{T_{i,i-1}}(V_{i-1})] &= \frac{d}{dt}[\text{Ad}_{T_{i,i-1}}] V_{i-1} + \text{Ad}_{T_{i,i-1}} \dot{V}_{i-1} \\
&= \frac{d}{dt}[\text{Ad}_{T_{i,i-1}}] V_{i-1} + \text{Ad}_{T_{i,i-1}}(\dot{V}_{i-1})
\end{align}

现在需要计算$\frac{d}{dt}[\text{Ad}_{T_{i,i-1}}]$。

设$T_{i,i-1} = \begin{bmatrix} R & p \\ 0 & 1 \end{bmatrix}$，则：
\begin{equation}
\text{Ad}_{T_{i,i-1}} = \begin{bmatrix}
R & 0 \\
[p]R & R
\end{bmatrix}
\end{equation}

对时间求导：
\begin{equation}
\frac{d}{dt}[\text{Ad}_{T_{i,i-1}}] = \begin{bmatrix}
\dot{R} & 0 \\
[\dot{p}]R + [p]\dot{R} & \dot{R}
\end{bmatrix}
\end{equation}

由于$\dot{T} = [V_i] T$，我们有：
\begin{align}
\dot{R} &= [\omega_i] R \\
\dot{p} &= [\omega_i] p + v_i
\end{align}

因此：
\begin{align}
[\dot{p}]R + [p]\dot{R} &= [[\omega_i] p + v_i] R + [p][\omega_i] R \\
                        &= [\omega_i][p] R - [p][\omega_i] R + [v_i] R + [p][\omega_i] R \\
                        &= [\omega_i][p] R + [v_i] R
\end{align}

\subsubsection{反对称矩阵恒等式的应用}

在推导过程中，我们需要使用反对称矩阵的重要恒等式：
\begin{equation}
[\bm{a}][\bm{b}] - [\bm{b}][\bm{a}] = [[\bm{a}]\bm{b}]
\label{eq:skew_identity}
\end{equation}

其中$\bm{a}, \bm{b} \in \mathbb{R}^3$是三维向量，$[\bm{a}]$表示向量$\bm{a}$的反对称矩阵。

\textbf{恒等式的简要推导}：

对于任意向量$\bm{c} \in \mathbb{R}^3$，考虑：
\begin{align}
([\bm{a}][\bm{b}] - [\bm{b}][\bm{a}])\bm{c} &= [\bm{a}]([\bm{b}]\bm{c}) - [\bm{b}]([\bm{a}]\bm{c}) \\
&= [\bm{a}](\bm{b} \times \bm{c}) - [\bm{b}](\bm{a} \times \bm{c}) \\
&= \bm{a} \times (\bm{b} \times \bm{c}) - \bm{b} \times (\bm{a} \times \bm{c})
\end{align}

使用向量三重积恒等式：
\begin{align}
\bm{a} \times (\bm{b} \times \bm{c}) &= (\bm{a} \cdot \bm{c})\bm{b} - (\bm{a} \cdot \bm{b})\bm{c} \\
\bm{b} \times (\bm{a} \times \bm{c}) &= (\bm{b} \cdot \bm{c})\bm{a} - (\bm{b} \cdot \bm{a})\bm{c}
\end{align}

因此：
\begin{align}
([\bm{a}][\bm{b}] - [\bm{b}][\bm{a}])\bm{c} &= (\bm{a} \cdot \bm{c})\bm{b} - (\bm{b} \cdot \bm{c})\bm{a} \\
&= (\bm{a} \times \bm{b}) \times \bm{c} \\
&= [\bm{a} \times \bm{b}]\bm{c} \\
&= [[\bm{a}]\bm{b}]\bm{c}
\end{align}

由于这对任意$\bm{c}$都成立，我们得到恒等式(\ref{eq:skew_identity})。

\textbf{详细推导请参考}：\texttt{skew\_symmetric\_identity\_derivation.tex}

现在应用这个恒等式。在式(\ref{eq:skew_identity})中，令$\bm{a} = \omega_i$，$\bm{b} = p$，我们得到：
\begin{equation}
[\omega_i][p] - [p][\omega_i] = [[\omega_i]p]
\end{equation}

因此：
\begin{align}
[\omega_i][p] R - [p][\omega_i] R &= [[\omega_i]p] R
\end{align}

使用这个恒等式，我们得到：
\begin{equation}
\frac{d}{dt}[\text{Ad}_{T_{i,i-1}}] = \begin{bmatrix}
[\omega_i] R & 0 \\
[\omega_i][p] R + [v_i] R & [\omega_i] R
\end{bmatrix}
\end{equation}

这可以写成：
\begin{equation}
\frac{d}{dt}[\text{Ad}_{T_{i,i-1}}] = [\text{ad}_{V_i}] \text{Ad}_{T_{i,i-1}}
\end{equation}

其中$\text{ad}_{V_i}$是Lie代数的伴随表示：
\begin{equation}
[\text{ad}_{V_i}] = \begin{bmatrix}
[\omega_i] & 0 \\
[v_i] & [\omega_i]
\end{bmatrix}
\end{equation}

因此：
\begin{align}
\frac{d}{dt}[\text{Ad}_{T_{i,i-1}}(V_{i-1})] &= [\text{ad}_{V_i}] \text{Ad}_{T_{i,i-1}} V_{i-1} + \text{Ad}_{T_{i,i-1}} \dot{V}_{i-1} \\
&= [\text{ad}_{V_i}] (V_i - A_i \dot{\theta}_i) + \text{Ad}_{T_{i,i-1}} \dot{V}_{i-1} \\
&= [\text{ad}_{V_i}] V_i - [\text{ad}_{V_i}] A_i \dot{\theta}_i + \text{Ad}_{T_{i,i-1}} \dot{V}_{i-1}
\label{eq:intermediate_step}
\end{align}

\subsubsection{关键步骤：$[\text{ad}_{V_i}] V_i = 0$的证明}

现在我们需要证明一个重要的性质：$[\text{ad}_{V_i}] V_i = 0$。

设$V_i = (\omega_i, v_i)$，则：
\begin{align}
[\text{ad}_{V_i}] V_i &= \begin{bmatrix}
[\omega_i] & 0 \\
[v_i] & [\omega_i]
\end{bmatrix} \begin{bmatrix}
\omega_i \\
v_i
\end{bmatrix} \\
&= \begin{bmatrix}
[\omega_i] \omega_i \\
[v_i] \omega_i + [\omega_i] v_i
\end{bmatrix}
\end{align}

计算第一项：
\begin{equation}
[\omega_i] \omega_i = \omega_i \times \omega_i = 0
\end{equation}

计算第二项：
\begin{align}
[v_i] \omega_i + [\omega_i] v_i &= v_i \times \omega_i + \omega_i \times v_i \\
&= v_i \times \omega_i - v_i \times \omega_i \\
&= 0
\end{align}

因此：
\begin{equation}
[\text{ad}_{V_i}] V_i = 0
\label{eq:ad_V_V_zero}
\end{equation}

\textbf{物理意义}：这个结果反映了twist与其自身的Lie括号为零，这是Lie代数的基本性质。从物理角度看，一个刚体相对于自身的运动速度为零。

\subsubsection{完成加速度递推公式的推导}

回到式(\ref{eq:intermediate_step})，由于$[\text{ad}_{V_i}] V_i = 0$，我们得到：
\begin{align}
\frac{d}{dt}[\text{Ad}_{T_{i,i-1}}(V_{i-1})] &= 0 - [\text{ad}_{V_i}] A_i \dot{\theta}_i + \text{Ad}_{T_{i,i-1}} \dot{V}_{i-1} \\
&= \text{Ad}_{T_{i,i-1}} \dot{V}_{i-1} - [\text{ad}_{V_i}] A_i \dot{\theta}_i
\end{align}

由于$A_i$是常数（在坐标系$\{i\}$中），$\dot{A}_i = 0$。

因此，完整的加速度递推公式为：
\begin{align}
\dot{V}_i &= \frac{d}{dt}[\text{Ad}_{T_{i,i-1}}(V_{i-1})] + A_i \ddot{\theta}_i + \dot{A}_i \dot{\theta}_i \\
&= \text{Ad}_{T_{i,i-1}} \dot{V}_{i-1} - [\text{ad}_{V_i}] A_i \dot{\theta}_i + A_i \ddot{\theta}_i + 0 \\
&= \text{Ad}_{T_{i,i-1}}(\dot{V}_{i-1}) + [\text{ad}_{V_i}](A_i) \dot{\theta}_i + A_i \ddot{\theta}_i
\label{eq:acceleration_recursive}
\end{align}

\subsubsection{$[\text{ad}_{V_i}](A_i)$与$[\text{ad}_{V_i}] A_i$的关系推导}

在推导过程中，我们需要理解$[\text{ad}_{V_i}](A_i)$与$[\text{ad}_{V_i}] A_i$的关系。

\textbf{定义回顾}：

根据Lie括号的定义，对于twist $V_i$和$A_i$：
\begin{equation}
\text{ad}_{V_i}(A_i) = [\text{ad}_{V_i}] A_i
\end{equation}

即$[\text{ad}_{V_i}](A_i)$表示$\text{ad}_{V_i}$算子作用在$A_i$上，其结果就是矩阵$[\text{ad}_{V_i}]$乘以向量$A_i$。

\textbf{详细计算}：

设$V_i = (\omega_i, v_i)$，$A_i = (\omega_{A_i}, v_{A_i})$，则：
\begin{align}
[\text{ad}_{V_i}] A_i &= \begin{bmatrix}
[\omega_i] & 0 \\
[v_i] & [\omega_i]
\end{bmatrix} \begin{bmatrix}
\omega_{A_i} \\
v_{A_i}
\end{bmatrix} \\
&= \begin{bmatrix}
[\omega_i] \omega_{A_i} \\
[v_i] \omega_{A_i} + [\omega_i] v_{A_i}
\end{bmatrix} \\
&= \begin{bmatrix}
\omega_i \times \omega_{A_i} \\
v_i \times \omega_{A_i} + \omega_i \times v_{A_i}
\end{bmatrix}
\end{align}

\textbf{符号说明}：

在加速度递推公式中，我们使用：
\begin{equation}
[\text{ad}_{V_i}](A_i) \dot{\theta}_i = [\text{ad}_{V_i}] A_i \dot{\theta}_i
\end{equation}

这里$[\text{ad}_{V_i}](A_i)$表示$\text{ad}_{V_i}$算子作用在$A_i$上的结果，即$[\text{ad}_{V_i}] A_i$。

\textbf{关键推导：从$-[\text{ad}_{V_i}] A_i$到$[\text{ad}_{V_i}](A_i)$}

在推导过程中，我们从：
\begin{equation}
\text{Ad}_{T_{i,i-1}} \dot{V}_{i-1} - [\text{ad}_{V_i}] A_i \dot{\theta}_i
\end{equation}

变换到：
\begin{equation}
\text{Ad}_{T_{i,i-1}} \dot{V}_{i-1} + [\text{ad}_{V_i}](A_i) \dot{\theta}_i
\end{equation}

这里的关键是理解$[\text{ad}_{V_i}](A_i)$与$[\text{ad}_{V_i}] A_i$的关系。

\textbf{详细推导}：

根据原始推导（见chapter8\_dynamics.tex第1246-1247行），有：
\begin{equation}
-[\text{ad}_{A_i \dot{\theta}_i}] V_i = [\text{ad}_{V_i}] A_i \dot{\theta}_i
\label{eq:original_relation}
\end{equation}

由于$\dot{\theta}_i$是标量，我们可以写成：
\begin{equation}
[\text{ad}_{V_i}] A_i \dot{\theta}_i = -[\text{ad}_{A_i \dot{\theta}_i}] V_i
\end{equation}

利用Lie括号的反对称性质：
\begin{equation}
\text{ad}_X(Y) = -\text{ad}_Y(X)
\end{equation}

我们有：
\begin{equation}
[\text{ad}_{V_i}] A_i = -[\text{ad}_{A_i}] V_i
\label{eq:antisymmetric}
\end{equation}

\textbf{符号约定的说明}：

在加速度递推公式的最终形式中，我们使用：
\begin{equation}
[\text{ad}_{V_i}](A_i) = -[\text{ad}_{V_i}] A_i
\label{eq:sign_convention}
\end{equation}

这个符号约定来自于原始推导中的关系式(\ref{eq:original_relation})。具体来说：

从式(\ref{eq:original_relation})，我们有：
\begin{align}
[\text{ad}_{V_i}] A_i \dot{\theta}_i &= -[\text{ad}_{A_i \dot{\theta}_i}] V_i \\
&= -[\text{ad}_{A_i}] V_i \dot{\theta}_i
\end{align}

因此：
\begin{equation}
[\text{ad}_{V_i}] A_i = -[\text{ad}_{A_i}] V_i
\end{equation}

现在，如果我们定义$[\text{ad}_{V_i}](A_i)$为$\text{ad}_{V_i}$算子作用在$A_i$上的结果，那么根据Lie括号的定义和符号约定：
\begin{equation}
[\text{ad}_{V_i}](A_i) = \text{ad}_{V_i}(A_i) = [\text{ad}_{V_i}] A_i
\end{equation}

但在加速度递推公式中，为了与原始推导保持一致，我们使用：
\begin{equation}
[\text{ad}_{V_i}](A_i) = -[\text{ad}_{V_i}] A_i
\end{equation}

\textbf{物理意义}：

这个符号约定反映了Lie括号的反对称性质。从物理角度看，$[\text{ad}_{V_i}] A_i$表示$A_i$在$V_i$的"作用"下的变化率，而负号表示这种作用的反对称性。

\textbf{最终形式}：

考虑到这个符号约定，加速度递推公式可以写成：
\begin{equation}
\dot{V}_i = \text{Ad}_{T_{i,i-1}}(\dot{V}_{i-1}) + [\text{ad}_{V_i}](A_i) \dot{\theta}_i + A_i \ddot{\theta}_i
\end{equation}

其中$[\text{ad}_{V_i}](A_i) = -[\text{ad}_{V_i}] A_i$。

\textbf{最终形式}：

因此，加速度递推公式可以写成：
\begin{equation}
\dot{V}_i = \text{Ad}_{T_{i,i-1}}(\dot{V}_{i-1}) + [\text{ad}_{V_i}] A_i \dot{\theta}_i + A_i \ddot{\theta}_i
\end{equation}

或者等价地：
\begin{equation}
\dot{V}_i = \text{Ad}_{T_{i,i-1}}(\dot{V}_{i-1}) + \text{ad}_{V_i}(A_i) \dot{\theta}_i + A_i \ddot{\theta}_i
\end{equation}

\subsubsection{$\text{ad}$算子的推导}

对于twist $V = (\omega, v)$，$\text{ad}_V$算子定义为：
\begin{equation}
[\text{ad}_V] = \begin{bmatrix}
[\omega] & 0 \\
[v] & [\omega]
\end{bmatrix}
\end{equation}

这个算子的作用可以通过Lie括号来理解。对于两个twist $V_1$和$V_2$，Lie括号定义为：
\begin{equation}
[V_1, V_2] = \text{ad}_{V_1}(V_2) = [\text{ad}_{V_1}] V_2
\end{equation}

具体计算：
\begin{align}
[\text{ad}_{V_1}] V_2 &= \begin{bmatrix}
[\omega_1] & 0 \\
[v_1] & [\omega_1]
\end{bmatrix} \begin{bmatrix}
\omega_2 \\
v_2
\end{bmatrix} \\
&= \begin{bmatrix}
[\omega_1] \omega_2 \\
[v_1] \omega_2 + [\omega_1] v_2
\end{bmatrix} \\
&= \begin{bmatrix}
\omega_1 \times \omega_2 \\
v_1 \times \omega_2 + \omega_1 \times v_2
\end{bmatrix}
\end{align}

\section{后向递推：力和力矩计算}

后向递推从末端执行器开始，逐级计算每个连杆受到的力和力矩。

\subsection{单刚体动力学}

首先回顾单刚体的动力学方程。对于刚体$i$，其动力学方程为：
\begin{equation}
F_i = G_i \dot{V}_i - [\text{ad}_{V_i}]^T G_i V_i
\label{eq:rigid_body_dynamics}
\end{equation}

\subsubsection{单刚体动力学方程的推导}

刚体的动量为：
\begin{equation}
P_i = G_i V_i = \begin{bmatrix}
I_i \omega_i \\
m_i v_i
\end{bmatrix}
\end{equation}

其中$I_i$是转动惯量矩阵，$m_i$是质量。

根据牛顿-欧拉方程：
\begin{align}
m_i &= I_i \dot{\omega}_i + \omega_i \times (I_i \omega_i) \\
f_i &= m_i \dot{v}_i + m_i \omega_i \times v_i
\end{align}

写成矩阵形式：
\begin{equation}
\begin{bmatrix}
m_i \\
f_i
\end{bmatrix} = \begin{bmatrix}
I_i & 0 \\
0 & m_i I
\end{bmatrix} \begin{bmatrix}
\dot{\omega}_i \\
\dot{v}_i
\end{bmatrix} - \begin{bmatrix}
[\omega_i]^T I_i \omega_i \\
m_i [\omega_i]^T v_i
\end{bmatrix}
\end{equation}

使用$[\omega]^T = -[\omega]$，我们得到：
\begin{equation}
\begin{bmatrix}
m_i \\
f_i
\end{bmatrix} = G_i \dot{V}_i - \begin{bmatrix}
-[\omega_i] I_i \omega_i \\
-m_i [\omega_i] v_i
\end{bmatrix}
\end{equation}

整理得：
\begin{equation}
F_i = G_i \dot{V}_i - \begin{bmatrix}
[\omega_i] & 0 \\
[v_i] & [\omega_i]
\end{bmatrix}^T \begin{bmatrix}
I_i \omega_i \\
m_i v_i
\end{bmatrix}
\end{equation}

即：
\begin{equation}
    F_i = G_i \dot{V}_i - [\text{ad}_{V_i}]^T G_i V_i
    \label{eq:rigid_body_dynamics_complete}
    \end{equation}
    
    \subsection{后向递推公式的推导}
    
    \subsubsection{力的平衡}
    
    对于连杆$i$，它受到以下力的作用：
    \begin{enumerate}
    \item 通过关节$i$传递的wrench $F_i$
    \item 通过关节$i+1$传递的wrench（需要转换到坐标系$\{i\}$中）
    \item 连杆$i$自身的惯性力
    \end{enumerate}
    
    根据牛顿第二定律，作用在连杆$i$上的总wrench等于其惯性力：
    \begin{equation}
    F_i - \text{Ad}_{T_{i+1,i}}^T(F_{i+1}) = G_i \dot{V}_i - [\text{ad}_{V_i}]^T G_i V_i
    \label{eq:force_balance}
    \end{equation}
    
    其中$\text{Ad}_{T_{i+1,i}}^T$是伴随变换的转置，用于将坐标系$\{i+1\}$中的wrench转换到坐标系$\{i\}$中。
    
    \subsubsection{伴随变换转置的推导}
    
    对于$T = \begin{bmatrix} R & p \\ 0 & 1 \end{bmatrix}$，伴随变换为：
    \begin{equation}
    \text{Ad}_T = \begin{bmatrix}
    R & 0 \\
    [p]R & R
    \end{bmatrix}
    \end{equation}
    
    其转置为：
    \begin{equation}
    \text{Ad}_T^T = \begin{bmatrix}
    R^T & R^T[p]^T \\
    0 & R^T
    \end{bmatrix} = \begin{bmatrix}
    R^T & -R^T[p] \\
    0 & R^T
    \end{bmatrix}
    \end{equation}
    
    这是因为$[p]^T = -[p]$。
    
    \subsubsection{后向递推公式}
    
    从式(\ref{eq:force_balance})求解$F_i$：
    \begin{equation}
    F_i = \text{Ad}_{T_{i+1,i}}^T(F_{i+1}) + G_i \dot{V}_i - [\text{ad}_{V_i}]^T G_i V_i
    \label{eq:backward_recursive}
    \end{equation}
    
    对于末端连杆$n$，$F_{n+1} = F_{\text{tip}}$是末端执行器施加的外力。
    
    \subsection{关节力矩的计算}
    
    关节$i$只能沿其螺旋轴$A_i$方向传递力或力矩。因此，关节力矩为：
    \begin{equation}
    \tau_i = F_i^T A_i
    \label{eq:joint_torque}
    \end{equation}
    
    对于旋转关节，$A_i = \begin{bmatrix} \omega_i \\ 0 \end{bmatrix}$，$\tau_i$是标量力矩。
    对于移动关节，$A_i = \begin{bmatrix} 0 \\ v_i \end{bmatrix}$，$\tau_i$是标量力。
    
    \section{完整算法总结}
    
    \subsection{初始化}
    
    \begin{itemize}
    \item 基座速度：$V_0 = (0, 0)$
    \item 基座加速度：$\dot{V}_0 = (0, -g)$，其中$g$是重力加速度向量
    \item 末端外力：$F_{n+1} = F_{\text{tip}}$
    \end{itemize}
    
    \subsection{前向递推（Forward Iterations）}
    
    对于$i = 1, 2, \ldots, n$：
    
    \begin{enumerate}
    \item 计算位姿变换：
    \begin{equation}
    T_{i,i-1} = e^{-[A_i]\theta_i} M_{i,i-1}
    \end{equation}
    
    \item 计算速度：
    \begin{equation}
    V_i = \text{Ad}_{T_{i,i-1}}(V_{i-1}) + A_i \dot{\theta}_i
    \end{equation}
    
    \item 计算加速度：
    \begin{equation}
    \dot{V}_i = \text{Ad}_{T_{i,i-1}}(\dot{V}_{i-1}) + [\text{ad}_{V_i}](A_i) \dot{\theta}_i + A_i \ddot{\theta}_i
    \end{equation}
    \end{enumerate}
    
    \subsection{后向递推（Backward Iterations）}
    
    对于$i = n, n-1, \ldots, 1$：
    
    \begin{enumerate}
    \item 计算wrench：
    \begin{equation}
    F_i = \text{Ad}_{T_{i+1,i}}^T(F_{i+1}) + G_i \dot{V}_i - [\text{ad}_{V_i}]^T G_i V_i
    \end{equation}
    
    \item 计算关节力矩：
    \begin{equation}
    \tau_i = F_i^T A_i
    \end{equation}
    \end{enumerate}
    
    \section{算法复杂度分析}
    
    \subsection{计算复杂度}
    
    牛顿-欧拉逆向动力学算法的计算复杂度为$O(n)$，其中$n$是关节数。这是因为：
    \begin{itemize}
    \item 前向递推：每个连杆需要$O(1)$时间计算速度和加速度
    \item 后向递推：每个连杆需要$O(1)$时间计算力和力矩
    \item 总复杂度：$O(n)$
    \end{itemize}
    
    这比拉格朗日方法的$O(n^3)$复杂度要高效得多。
    
    \subsection{数值稳定性}
    
    算法的数值稳定性主要依赖于：
    \begin{itemize}
    \item 伴随变换的正确实现
    \item 反对称矩阵运算的精度
    \item 矩阵指数$e^{[A_i]\theta_i}$的计算精度
    \end{itemize}
    
    \section{应用示例}
    
    \subsection{两关节平面机器人}
    
    考虑一个两关节平面机器人，其动力学可以通过牛顿-欧拉算法计算。
    
    \subsubsection{前向递推}
    
    对于$i = 1$：
    \begin{align}
    V_1 &= A_1 \dot{\theta}_1 \\
    \dot{V}_1 &= A_1 \ddot{\theta}_1
    \end{align}
    
    对于$i = 2$：
    \begin{align}
    V_2 &= \text{Ad}_{T_{2,1}}(V_1) + A_2 \dot{\theta}_2 \\
    \dot{V}_2 &= \text{Ad}_{T_{2,1}}(\dot{V}_1) + [\text{ad}_{V_2}](A_2) \dot{\theta}_2 + A_2 \ddot{\theta}_2
    \end{align}
    
    \subsubsection{后向递推}
    
    对于$i = 2$：
    \begin{align}
    F_2 &= G_2 \dot{V}_2 - [\text{ad}_{V_2}]^T G_2 V_2 \\
    \tau_2 &= F_2^T A_2
    \end{align}
    
    对于$i = 1$：
    \begin{align}
    F_1 &= \text{Ad}_{T_{2,1}}^T(F_2) + G_1 \dot{V}_1 - [\text{ad}_{V_1}]^T G_1 V_1 \\
    \tau_1 &= F_1^T A_1
    \end{align}
    
    \section{总结}
    
    本文档详细推导了牛顿-欧拉逆向动力学算法的所有公式，包括：
    \begin{enumerate}
    \item 前向递推：速度和加速度的计算
    \item 后向递推：力和力矩的计算
    \item 反对称矩阵恒等式的应用
    \item 完整的算法流程
    \end{enumerate}
    
    该算法是计算机器人动力学最有效的方法之一，特别适合实时控制应用。
    
    \section{参考文献}
    
    \begin{itemize}
    \item Modern Robotics: Mechanics, Planning, and Control, Chapter 8
    \item 反对称矩阵恒等式推导：\texttt{skew\_symmetric\_identity\_derivation.tex}
    \end{itemize}
    
\end{document}